\section{Conclusion}
\label{sec:conclusion}

We have presented Morphogenic Intelligence, a neural architecture that treats network topology as a learned artifact rather than a fixed design choice. Morphons are born, differentiate, fuse, and die under local rules grounded in six biological principles. The full system is implemented in Rust, runs in real time on commodity hardware, and is reproducible from a clean checkout.

We solved CartPole-v1 (avg=195.2 on the v0.5.0 codebase) using only local three-factor learning and developmental morphogenesis, with no backpropagation through the morphon network. On MNIST we reach \textbf{44.5\%} intact accuracy in v4.1.0 --- up from 30\% in v3.0.0 --- driven by two complementary improvements: spike-timing inhibitory STDP (Vogels 2011) that breaks hub-morphon dominance (1 labeled class $\rightarrow$ 12), and Pulse Kernel Lite that delivers a \textbf{4.4$\times$ wall-clock speedup} by fitting the fast-path arrays in L1 cache. On a synthetic NLP readiness benchmark, the morphon substrate produces representations that an analog readout can exploit at 99\% / 62\% / 85\% across three tiers, even though the spike-based output pathway loses the same information.

We documented four failure modes encountered during development. Three (modulatory explosion, motor silencing, LTD vicious cycle) have well-known biological analogues and biology-derived fixes. The fourth --- premature ``Mature'' stage detection on dense-reward classification --- is, to our knowledge, novel and broadly applicable to any biologically-inspired homeostatic system that uses raw reward as a maturity signal. The fix is to use reward prediction error instead, which is the principled approach all such systems should adopt. A fifth failure mode, identified in v4.1.0, is the \emph{timescale mismatch} in inhibitory STDP: computing inhibitory weight updates on a slow running average (100-tick window, medium period) allows hub morphons to entrench before corrective inhibition can respond. The fix --- per-spike inhibitory plasticity at biological timescales --- is both the principled and the practical solution.

The most surprising finding is that beneficial self-healing exists. Under the hub-breaking iSTDP regime, a post-damage network (30\% ablated) reaches 48.0\% --- above the intact baseline --- because ablation removes entrenched hubs and Endoquilibrium restores the high-plasticity regime that the original training never sustained. This points to a more general principle: \emph{the optimal training trajectory for a developmental system may not be a smooth ascent but a rhythm of structural shocks and recovery}. Whether this principle extends to the architectures that currently dominate machine learning is an open question we cannot answer, but it is one we think the field should ask.

The MI engine, the experiments in this paper, and the full failure-mode catalogue are available at \texttt{https://github.com/SimplyLiz/Morphon}. This preprint is archived at Zenodo with DOI \texttt{10.5281/zenodo.19467243}.

\paragraph{Acknowledgments.}
The authors thank the Rust ecosystem, the petgraph crate, and the Diehl \& Cook MNIST baseline that calibrated our expectations.
